% ======================================================================
% SECTION 5: LIMITATIONS AND FUTURE WORK
% Five thematic subsections: approximations vs generated vs executed
% accounting, 60min vs 1min formal deltas, cross-simulation
% limitations, Track A vs Track B future directions, future work to
% reduce approximations.
% Location: 2030-gbm-1min/paper/full-paper/sections/limitations_future.tex
% ======================================================================

\subsection{Approximations vs Generated vs Executed Accounting}
\label{sec:limits-accounting}

Table \ref{tab:accounting} renders the three-way accounting for
every artifact category in the 2030-gbm-1min/ tree. The columns
distinguish what was approximated by the upstream instruction set,
what was generated by Claude Code Opus 4.7 1M Max
\cite{claude-opus-47, claude-code, one-minute-variant-instructions},
and what was executed by the simulation pipeline.

\begin{table}[H]
\centering
\caption{Approximated vs generated vs executed artifact accounting
across the 2030-gbm-1min/ tree.}
\label{tab:accounting}
\begin{tabular}{>{\raggedright\arraybackslash}p{3.0cm}
                >{\raggedright\arraybackslash}p{4.0cm}
                >{\raggedright\arraybackslash}p{4.3cm}
                >{\raggedright\arraybackslash}p{3.5cm}}
\toprule
\textbf{Category} & \textbf{Approximated upstream} & \textbf{Generated by Claude Code} & \textbf{Executed in outputs/} \\
\midrule
Instruction files & 12 Markdown files & na & na \\
Toolchain & file_size_pyramid_1min.md (caps) & pyproject.toml, docker-compose.yml, Cargo.toml & outputs/logs/ \\
Sensor schema (200 channels) & commit_02_sensors_1min.md & schemas/sensor_record_4arm.\{schema.json, proto, avsc\} & outputs/sensors/ \\
7-DOF DH parameter table & robot_specification_neurospeed.md & config/kinematics_4arm.yaml & outputs/xyz_mapping/ \\
16-iteration sweep design & commit_04_iterations_1min.md & config/iterations.yaml, src/simulation/iterate_1min.py & outputs/iterations/ \\
Composite score formula & commit_05_competition_1min.md (frozen weights) & src/metrics/compute_1min.py & outputs/metrics/ \\
LLM tournament & commit_05_competition_1min.md & src/llm/compare_agent_1min.py + prompts/comparison_prompt_1min.md & outputs/comparison/ + outputs/comparison_robot_vs_human/ \\
Zenodo L0 deposition & zenodo_archive_protocol.md (416 MB at DOI 10.5281/zenodo.18445179) & src/zenodo/patch_pointers.py + pointer JSON files & outputs/iterations/*_L0_raw.zenodo_pointer.json \\
Sensor sample table & sensor_specification_10khz.md (schema only) & src/sensors/ingest_4arm.py & outputs/sensors/sensor_sample_4arm.csv (54 by 1001) \\
ASCII diagrams & multi_arm_coordination.md (heartbeat sketch) & docs/architecture_overview_4arm.txt & 6 ASCII at outputs/diagrams/ \\
\bottomrule
\end{tabular}
\end{table}

Every place where the upstream instruction approximates a count or a
size, the downstream pass committed to a specific number; the
approximations are flagged for future tightening in
\S\ref{sec:limits-future}.

% \clearpage

\subsection{60-Minute vs 1-Minute Formal Delta Accounting}
\label{sec:limits-deltas}

The present 1-minute scope extends the prior 1-hour scope
\cite{kawchak_2026_19994945, repo-robotic-surgeries} along the
dimensions listed in Table \ref{tab:deltas}. This subsection is the
formal home of the 60min vs 1min delta accounting; downstream prose
elsewhere should cross reference here.

\begin{table}[H]
\centering
\caption{Formal delta between the prior 60-minute scope and the
current 1-minute scope.}
\label{tab:deltas}
\begin{tabular}{>{\raggedright\arraybackslash}p{3.9cm}
                >{\raggedright\arraybackslash}p{4.4cm}
                >{\raggedright\arraybackslash}p{4.4cm}
                >{\raggedright\arraybackslash}p{2.1cm}}
\toprule
\textbf{Dimension} & \textbf{60 min scope} & \textbf{1 min scope} & \textbf{Ratio} \\
\midrule
Cooperating arms & 1 & 4 & 4x \\
Tool assignment per arm & shared & arm 1 hyb u-w-p, arm 2 bipolar, arm 3 suction, arm 4 imaging & dedicated \\
Sensor channels & 50 (single arm, 1 kHz) & 200 across 4 arms (1 kHz cmd + 10 kHz force) & 4x \\
Force sample rate & 1 kHz & 10 kHz force & 10x \\
Cartesian update rate & 1 kHz cmd & 1 kHz cmd + 10 kHz force & finer \\
Iteration count & 64 & 16 & 0.25x \\
Tournament entities & 4 sponsors & 4 robots (mixed allowed) & same \\
Structural caveat & pharmacology dim only & time dim 1 min vs 1 h & new \\
Procedure duration & 60 minutes & 60 seconds & 60x faster \\
Per-arm tip force cap & 15.0 N & 5.0 N & 3x tighter \\
Cumulative cross-arm cap & not specified & 12.0 N envelope & new floor \\
E-stop latency & 50 ms & 5 ms & 10x faster \\
Committed footprint & multi-thousand artifact set & about 8.2 MB + 1.5 MB overhead & smaller \\
\bottomrule
\end{tabular}
\end{table}

\subsection{Cross-Simulation Limitations}
\label{sec:limits-cross}

The limitations that apply across every artifact in the
2030-gbm-1min/ tree are: (a) the synthetic patient PAT-GBM-0001
(4.2 cm right frontal IDH-wildtype glioblastoma) is not a real
patient; no IRB, no informed consent, and no real EHR signal was
used. (b) Claude Code Opus 4.7 1M Max output is not deterministic
across runs; a re-generation of the 2030-gbm-1min/ tree may produce
different file names, sizes, and prose, although the per-iteration
Parquet aggregates under seed 20260510 remain bit deterministic for
fixed parameter tuple \cite{claude-opus-47, claude-code}. (c) The
synthetic 16-iteration sweep has not been validated against any
real-patient cohort or the TRIPOD+AI / CREMLS reporting floor
\cite{Collins2024TRIPODAI, ElEmam2024CREMLS}. 

(d) The hypothetical
2030 Medtronic NeuroSpeed 1.0 hardware does not exist; the
comparison against Medtronic ROSA ONE Brain v3.0
\cite{rosa-one-brain, Lefranc2014ROSAStereotaxy} is paper-only.
(e) The C++20 control loop, C++20 1 kHz heartbeat layer, and Rust
2021 high-throughput runner are documented but may not have been
exhaustively compiled across every supported platform in this
generation pass; the partial-verification status is reproduced from
2030-gbm-1min/outputs/reports/limitations.md. (f) The LLM tournament
rationale at outputs/comparison/comparison.json and
outputs/comparison_robot_vs_human/comparison.json carries the
structural-time-dimension caveat in every round and the caveat must
not be removed downstream. (g) The work is not endorsed or sponsored
by any trial sponsor, FDA, CRO, site, IRB, regulator, or medical
society and was generated using Artificial Intelligence.

% \clearpage

\subsection{Track A vs Track B Future Directions}
\label{sec:limits-tracks}

Two future-work tracks extend the computational advantages
demonstrated here into real-world applications.

\textbf{Track A: single big model performs all tasks.} A future
generation of Claude Code (or an equivalent foundation-model agent)
holds the entire 2030-gbm-1min/ repository in context and emits all
per-iteration patient predictions, per-iteration safety signals, and
per-iteration FDA-bound signals from a single 1M token inference
pass. Architecturally simple but with higher inference cost per
decision and a single point of regulatory accountability.

\textbf{Track B: single big model orchestrates smaller agents.} A
future generation uses the 1M token context model only for
orchestration and routing, while delegating per-task prediction
(per-arm force prediction, per-iteration composite score, per-round
tournament verdict) to smaller specialized agents that run on
commodity hardware (Intel Core i5-6200U class or similar).
Distributes regulatory accountability but lowers per-site inference
cost and improves the security posture for PHI handling.

\begin{table}[H]
\centering
\caption{Track A vs Track B trade-offs. Property column on the left;
Track A and Track B side by side.}
\label{tab:tracks}
\begin{tabular}{>{\raggedright\arraybackslash}p{2.6cm}
                >{\raggedright\arraybackslash}p{6.3cm}
                >{\raggedright\arraybackslash}p{6.3cm}}
\toprule
\textbf{Property} & \textbf{Track A: big single model} & \textbf{Track B: big model + small agents} \\
\midrule
Architecture & one 1M token pass holds everything & router plus N specialized agents \\
Inference cost & higher per decision & lower per decision \\
Regulatory accountability & single SaMD entry point & multiple SaMD entry points \\
PHI handling & needs on-prem or PHI-cleared cloud & easier on-prem per-task isolation \\
Hardware floor & one 1M token capable backend & small backends per task \\
Update cycle & one re-deploy & per-agent re-deploys \\
Extension target & cloud-orchestrated NeuroSpeed-class & on-site agent network on existing CI \\
\bottomrule
\end{tabular}
\end{table}

The local-backend options for Track B include Ollama \cite{ollama}
and vLLM \cite{vllm} for the small agents, with Claude Sonnet 4.6
\cite{claude-sonnet-46} or Claude Opus 4.7 \cite{claude-opus-47} as
the orchestrator.

\subsection{Future Work to Reduce Approximations}
\label{sec:limits-future}

The future work plan reduces the approximations enumerated in
\S\ref{sec:limits-accounting}, the 60-minute vs 1-minute deltas in
\S\ref{sec:limits-deltas}, the cross-simulation limitations in
\S\ref{sec:limits-cross}, and the Track A vs Track B trade-offs in
\S\ref{sec:limits-tracks}. The concrete deliverables are:

(a) Extend the 16-iteration sweep at outputs/iterations/ to 32 or
64 iterations to tighten cumulative-arm-force violation accounting.

(b) Add a competing 4-arm robot vendor entry to the LLM tournament
at outputs/comparison/ so the leaderboard is not single vendor.

(c) Run an ablation study on the cumulative-force-share clamp (soft
11.0 N threshold) vs the hard 12 N park threshold at
outputs/iterations/.

(d) Increase the sensor sample rate beyond 10 kHz force toward 100
kHz force per arm to capture the high-frequency tip dynamics during
Phase 2 bulk resection.

(e) Compress the 60-second target to 30 seconds and 15 seconds in
follow-on variants that share this repository under 2030-gbm-30s/
and 2030-gbm-15s/.

% \clearpage

(f) Ingest the FDA RTCT real-time signal-sharing channel
\cite{fda2026realtime} and emit per-iteration endpoints at hourly
cadence.

(g) Adopt TRIPOD+AI and CREMLS reporting standards
\cite{Collins2024TRIPODAI, ElEmam2024CREMLS} as the floor for any
real-patient extension.

(h) Deposit the release-aggregate L0 raw archive (currently 416 MB
across 16 iterations) on Zenodo \cite{zenodo} at DOI
10.5281/zenodo.18445179 with per-iteration pointer JSON resolution.

(i) Replace the synthetic patient PAT-GBM-0001 with a TRIPOD+AI-
compliant retrospective cohort that pairs the 16-iteration sweep
against real EHR-linked oncology outcomes.

(j) Prototype both Track A and Track B on the on-premises stack and
report the per-iteration cost and latency delta.
